%% Chapter 4, Section 4.4: Discussion and Bibliography

\section{Discussion and Bibliography}
\label{sec:4.4}

The Metropolis Hastings algorithm was introduced in~\cite{162} restricted to the case of Gaussian
random walk proposals; the extension to user-chosen proposal kernel was proposed in~\cite{108}.
The paper~\cite{100} introduced an adaptive strategy, where the proposal kernel is updated along
the process using the full information gathered so far. Historical overviews of the Metropolis
Hastings algorithm and related MCMC methods include~\cite{192,61}. Many books~\cite{88,37,193}
and survey articles~\cite{36,225} are devoted to the Metropolis Hastings algorithm; adaptive strategies
are surveyed in~\cite{13,200}.

As shown in Theorem~\ref{thm:4.3}, the Metropolis Hastings algorithm is designed to ensure detailed
balance with respect to the target; in other words, Markov chains generated using the Metropolis
Hastings algorithm are \textit{reversible}. The paper~\cite{26} shows that the algorithm can be modified to
generate \textit{non-reversible} chains that may have better properties in terms of mixing behavior or
asymptotic variance. Uniform ergodicity of the independence sampler and geometric ergodicity
of RWMH were established in~\cite{161}. A comparison between rejection sampling, importance
sampling, and the independence sampler can be found in~\cite{144}. The study of optimal scaling
for Metropolis Hastings algorithms, which originates from~\cite{196,203} and is reviewed in~\cite{198},
is still an active area of research~\cite{240}. The convergence diagnostics by Raftery, Geweke, and
Gelman and Rubin were introduced in~\cite{189,84,81}. We refer to~\cite{167,241,230,229} for other
important examples of convergence diagnostics, and to~\cite{38} for an early survey on this topic.

An important feature of the Metropolis Hastings algorithm is that it can be implemented
without knowledge of the normalizing constant of the target. However, in many applications
such as Bayesian inverse problems~\cite{219}, machine learning with tall data~\cite{16}, and genetics~\cite{17}
evaluating the unnormalized target density to compute the Metropolis Hastings acceptance
probability can be expensive. If a \textit{deterministic} cheap-to-evaluate approximation of the target is
available, this issue can be alleviated using delayed-acceptance Metropolis Hastings algorithms~\cite{50,69,68,212}. The idea is similar to that behind the envelope rejection sampling algorithm
discussed in Chapter~2: first, we do an accept/reject step with the cheap-to-evaluate target approximation; if the move is accepted, we do a second accept/reject step with the expensive-to-evaluate
target. If a \textit{stochastic} unbiased estimator of the unnormalized target density is available, pseudomarginal MCMC algorithms~\cite{17,12} can be employed. Particular instances of pseudo-marginal
MCMC algorithms are particle MCMC algorithms~\cite{10} that rely on unbiased estimates computed via particle filters (see Chapter~9). Particle MCMC algorithms are powerful algorithms for
parameter estimation in hidden Markov models, see e.g.~\cite{93}.

In this chapter, we introduced the Metropolis Hastings algorithm on a countable or uncountable state space $E \subset \R^d$. Reversible jump MCMC algorithms~\cite{96} allow to sample posterior
distributions on spaces of varying dimension; this important extension of the Metropolis Hastings framework enables Bayesian inference in applications where the number of parameters in
the model is unknown. An adaptive, delayed-rejection algorithm for reversible jump MCMC was
introduced in~\cite{97}. Another important extension of the Metropolis Hastings framework concerns
its formulation in general state spaces~\cite{226,199}, which enables for instance function space sampling~\cite{52}. A common limitation of MCMC methodology is that it is not easy to parallelize due
to its serial structure. We refer to~\cite{41,91,211} for attempts to parallelize Metropolis Hastings
algorithms. These approaches often propose several moves in each iteration, an idea also found
in~\cite{147}. Finally, we remark that MCMC algorithms that rely on piecewise deterministic Markov
processes are studied for instance in~\cite{74,25,32}.
